% ==============================================================================
% header-common.tex
% Gemeinsame Pakete für Vortrag und Ausarbeitung
% Common packages for presentation and paper
%
% ARCHITEKTUR-HINWEIS / ARCHITECTURE NOTE:
%   Diese Datei enthält ausschließlich \usepackage-Aufrufe, die sowohl mit dem
%   Beamer-Dokumentenstil (Vortrag/) als auch mit scrartcl (Ausarbeitung/) sicher
%   kompatibel sind.
%
%   Konfigurationsbefehle (\definecolor, \hypersetup, \setlist, \newtheorem,
%   \tikzset, \newcommand, \lstset usw.) dürfen NICHT hier stehen, weil Beamer
%   viele davon intern bereits setzt; doppelte Definitionen erzeugen Konflikte
%   oder überschreiben das Beamer-Theme stillschweigend.
%   → Solche Befehle gehören in header-article.tex (Ausarbeitung) bzw. in das
%     Beamer-Theme (Vortrag).
%
% BEWUSST NICHT IN DIESER DATEI / INTENTIONALLY ABSENT:
%
%   inputenc / fontenc
%       Engine-spezifisch: LuaLaTeX verarbeitet UTF-8 nativ und benötigt diese
%       Pakete nicht (fontspec übernimmt die Aufgabe). pdfLaTeX braucht sie.
%       → Werden per \ifluatex-Konditional in Ausarbeitung.tex / Vortrag.tex
%         gesetzt, damit die Engine frei wählbar bleibt.
%
%   babel
%       Die Dokumentsprache ist projektspezifisch (ngerman, english, …) und
%       gehört daher ins Hauptdokument, wo sie für Nutzer leicht zu finden und
%       zu ändern ist.
%       → Ausarbeitung.tex / Vortrag.tex mit erklärendem Kommentar.
%
%   amsthm
%       Beamer definiert intern eine eigene Theorem-Infrastruktur. Wird amsthm
%       davor geladen, überschreibt es diese, was zu \newtheorem-Kollisionen und
%       stilistischem Bruch mit dem Beamer-Theme führt.
%       → Nur in header-article.tex (Ausarbeitung) geladen.
%
%   tipa
%       Stellt phonetische IPA-Symbole bereit, definiert dafür aber viele
%       Textbefehle um (\textprimstress, \textscripta usw.) und ist bekannt für
%       Konflikte mit lmodern und fontspec. Das Template verwendet tipa an keiner
%       Stelle – das Paket wurde daher vollständig entfernt.
%
%   natbib / cite / \usepackage{bibtex}
%       Alle drei waren bereits auskommentierter Altcode und sind inkompatibel
%       mit biblatex (dem aktuellen Literaturpaket). „bibtex" ist zudem kein
%       LaTeX-Paket, sondern ein externes Programm.
%       → Entfernt; biblatex + biber ersetzen den gesamten bibtex-Workflow.
%
% ==============================================================================

% ──────────────────────────────────────────────────────────────────────────────
% Engine-Erkennung
% ──────────────────────────────────────────────────────────────────────────────

% Stellt \ifluatex, \ifpdftex usw. bereit, mit denen engine-spezifische Pakete
% bedingt geladen werden können (z.B. inconsolata und lmodern nur unter pdfLaTeX).
\usepackage{iftex}

% ──────────────────────────────────────────────────────────────────────────────
% Schriften / Fonts
% ──────────────────────────────────────────────────────────────────────────────

% Inconsolata: gut lesbare Festbreitenschrift für Code-Listings.
% Wird unter LuaLaTeX übersprungen, weil dort fontspec + \setmonofont verwendet
% werden (gesetzt in Ausarbeitung.tex / Vortrag.tex).
\ifluatex\else
  \usepackage{inconsolata}
\fi

% Latin Modern: Erweiterter Zeichenvorrat gegenüber Computer Modern, korrekte
% Silbentrennung für Sonderzeichen. Wird unter LuaLaTeX übersprungen (fontspec
% übernimmt die Aufgabe vollständig).
\ifluatex\else
  \usepackage{lmodern}
\fi

% ──────────────────────────────────────────────────────────────────────────────
% Unicode und Sonderzeichen
% ──────────────────────────────────────────────────────────────────────────────

% Erlaubt die Definition benutzerdefinierter Unicode-Zeichen-Mappings mit
% \newunicodechar{<Zeichen>}{<LaTeX-Befehl>}.
% Wird in header-article.tex z.B. für den schmalen Leerabstand U+202F genutzt.
\usepackage{newunicodechar}

% ──────────────────────────────────────────────────────────────────────────────
% Grafiken und Abbildungen
% ──────────────────────────────────────────────────────────────────────────────

% Einbinden externer Bilddateien (\includegraphics, \resizebox usw.)
\usepackage{graphicx}

% Erweiterung des picture-Environments: ermöglicht Kreise, Kurven und Linien
% mit beliebiger Steigung ohne externe Pakete.
\usepackage{pict2e}

% TikZ: mächtige Beschreibungssprache für Vektorgrafiken direkt in LaTeX
\usepackage{tikz}

% TikZ-Bibliothek: verkettete Listen, Kettendiagramme und mehrteilige Knoten
\usetikzlibrary{chains,shapes.multipart}

% TikZ-Bibliothek: Flowchart-Formen (Rauten, Ellipsen), Pfeilspitzen, relative
% Positionierung von Knoten
\usetikzlibrary{shapes,arrows,positioning}

% PGFPlots: Datenvisualisierung (Funktions- und Datengraphen) auf TikZ-Basis
\usepackage{pgfplots}
% compat=1.18 stellt sicher, dass neue Achsenskalierung und Tick-Formatierung
% aktiv sind; verhindert Kompatibilitätswarnungen bei neueren TeX-Installationen.
\pgfplotsset{compat=1.18}

% Mehrere Unterabbildungen nebeneinander mit individuellen Beschriftungen
% (\subcaption, \subfigure-Umgebung); kompatibel mit Beamer.
\usepackage{subcaption}

% ──────────────────────────────────────────────────────────────────────────────
% Mathematik
% ──────────────────────────────────────────────────────────────────────────────

% Kern-AMS-Mathematik: equation, align, gather, cases, Operatoren, \text, …
% Muss vor mathtools geladen werden.
\usepackage{amsmath}

% AMS-Mathematikschriften: \mathbb (Doppelstrich), \mathfrak (Fraktur), …
\usepackage{amsfonts}

% Zusätzliche AMS-Symbole: \leqslant, \varnothing, \mathbb{R}, Pfeile, …
\usepackage{amssymb}

% Doppelstrich-Schrift für Mengensymbole: \mathds{N}, \mathds{1} (Indikatorfunktion)
% Ergänzt \mathbb aus amsfonts um serifenlose Varianten.
\usepackage{dsfont}

% Erweiterungen zu amsmath: \coloneqq, \shortintertext, Matrizenabstände,
% centerdots-Option sorgt für zentrierte Multiplikationspunkte (·).
\usepackage[centerdots]{mathtools}

% Schöne Inline-Brüche im Textsatz: \nicefrac{1}{2} → ½-Darstellung
\usepackage{nicefrac}

% St-Mary-Road-Symbole: Semantik-Klammern ⟦⟧ (\llbracket, \rrbracket),
% Blitz-Pfeil, weitere Logik-/Typentheorie-Symbole
\usepackage{stmaryrd}

% ──────────────────────────────────────────────────────────────────────────────
% Textsymbole
% ──────────────────────────────────────────────────────────────────────────────

% Klassische LaTeX-Symbole: \lhd, \rhd, \Box, \Diamond u.a., die weder in
% amssymb noch in base-LaTeX definiert sind.
\usepackage{latexsym}

% Textsymbole der TS1-Kodierung: \textbullet, \texttimes, \textdegree, Pfeile
% im Textmodus. Wird von vielen anderen Paketen implizit vorausgesetzt.
\usepackage{textcomp}

% ──────────────────────────────────────────────────────────────────────────────
% Farben und Tabellen
% ──────────────────────────────────────────────────────────────────────────────

% Farbnamen, Farbmischung und Farbräume (rgb, cmyk, HTML-Hex).
% Lädt das Paket colortbl automatisch mit, daher steht colortbl danach.
\usepackage{xcolor}

% Farbige Tabellenzellen und -zeilen (\cellcolor, \rowcolor, \columncolor)
\usepackage{colortbl}

% Professionelle Tabellenlinien: \toprule, \midrule, \bottomrule (kein \hline)
\usepackage{booktabs}

% Erweiterte Spaltentypen für tabular: >{…}, <{…}, m{}, p{}, b{}
\usepackage{array}

% Float-Positionierung: H-Positionsparameter (genau hier, kein Verschieben)
\usepackage{float}

% ──────────────────────────────────────────────────────────────────────────────
% Code-Listings
% ──────────────────────────────────────────────────────────────────────────────

% Syntax-hervorgehobene Quelltextdarstellung; Konfiguration (\lstset,
% \lstdefinelanguage) steht in header-article.tex, da Beamer eigene
% verbatim-Blöcke verwendet und die globale lstset-Konfiguration dort stört.
\usepackage{listings}

% ──────────────────────────────────────────────────────────────────────────────
% Algorithmen in Pseudocode
% ──────────────────────────────────────────────────────────────────────────────

% algorithm: Float-Umgebung für Algorithmen (wie figure/table, mit Beschriftung)
\usepackage{algorithm}

% algpseudocode: konkrete Pseudocode-Syntax auf Basis des algorithmicx-Pakets
% (\Procedure, \If, \For, \State, \Return usw.)
\usepackage{algpseudocode}

% ──────────────────────────────────────────────────────────────────────────────
% Listen und Aufzählungen
% ──────────────────────────────────────────────────────────────────────────────

% Vollständige Kontrolle über itemize/enumerate/description: Abstände, Labels,
% Einzüge. Konkrete Stilvorgaben (\setlist) stehen in header-article.tex,
% da Beamer Listenabstände intern für Folien optimiert.
\usepackage{enumitem}

% ──────────────────────────────────────────────────────────────────────────────
% Textboxen und Hilfsmittel
% ──────────────────────────────────────────────────────────────────────────────

% Farbige Rahmenboxen mit Titel, Hintergrundfarbe und vielen Stilen
% (tcolorbox, tcblisting usw.); kompatibel mit Beamer-Blöcken.
\usepackage{tcolorbox}

% Blindtext für Layout-Tests: \blindtext, \Blindtext, \blinddocument
% Kann nach Fertigstellung der Ausarbeitung entfernt werden.
\usepackage{blindtext}

% Randnotizen und farbige Inline-Hinweise während des Schreibens:
% \todo{...}, \missingfigure{...}. Pakete dieser Art werden typischerweise
% vor der Abgabe entfernt oder per \usepackage[disable]{todonotes} deaktiviert.
\usepackage{todonotes}

% ──────────────────────────────────────────────────────────────────────────────
% Bibliographie und Hyperlinks
% ──────────────────────────────────────────────────────────────────────────────

% Korrekte kontextsensitive Anführungszeichen (\enquote{...}); von biblatex
% explizit als Voraussetzung gefordert – ohne csquotes gibt biblatex eine
% Warnung aus.
\usepackage{csquotes}

% Klickbare Verweise, Sprünge im PDF und PDF-Metadaten.
% Muss möglichst spät geladen werden (nach fast allen anderen Paketen), damit
% die Link-Makros korrekt greifen. Farbkonfiguration (\hypersetup) folgt in
% header-article.tex, da Beamer die Link-Farben über das Theme steuert.
\usepackage{hyperref}
